\section{The zeta identities with error bounds}\label{sec:analytic}

The analytic variable can be kept outside the language of iPRA. A rational
complex approximation is a pair of rational numbers. A name for a number
supplies primitive recursive approximants within $2^{-j}$ at precision $j$.
A statement about arbitrary such inputs is a schema in these names, or a
finite calculation with input approximations and their error bounds. In
particular, no evaluator for arbitrary primitive recursive programs is
needed inside arithmetic. Equality of named quantities means agreement at
every precision, with errors included. Throughout this section assume a
supplied integer $q\ge1$ with $\sigma=\operatorname{Re}s\ge1+1/q$
and a supplied integer upper bound on $|s|$. Both are quantitative
inputs to the estimates.

\paragraph{Elementary analytic operations.}
Here are details of the rational interpretation used below. For positive
rational $u$, write $u=2^jv$, $1\le v<2$, by bounded doubling of the
numerator or denominator. Compute $\log u=j\log 2+\log v$ using
\[
   \log v=2\sum_{h=0}^{K-1}\frac{w^{2h+1}}{2h+1}+E_K,
   \qquad w=\frac{v-1}{v+1},\qquad
   |E_K|\le \frac{2|w|^{2K+1}}{1-|w|^2}.
\]
Use the same formula at $v=2$. It follows by summing the finite geometric
identity for $1/(1-w^2)$ and integrating. For $|z|\le B$ with integral
$B\ge1$, the exponential Taylor series beyond index $2B$ has successive
absolute terms in ratio at most $1/2$. The term at $2B$ is at most
$B^{2B}$, so another $2B^2+j+4$ terms give error below $2^{-j}$.
Thus $x^{-s}=\exp(-s\log x)$ for positive $x$ has an explicit
rational algorithm, also for named $x$ on an interval bounded away from
zero. The formulae $(\log x)'=1/x$, $(e^z)'=e^z$, the exponential
addition law, and $|e^{u+iv}|=e^u$ follow by the same geometric
remainders for the differentiated series, the binomial formula for their
product, and $e^{iv}e^{-iv}=1$. For real arguments one may equivalently
define $\log$ by integration of $1/x$. Differentiating
$e^{\log x}/x$ shows it is constant and evaluation at $1$ gives its
value. In these uses the fundamental theorem of calculus has a finite
proof by subdivision. On an interval of length $L$ where $|F''|\le B$, summing
the first-order increments of $F$ over a mesh of width $h$ gives
a left sum for $F'$ with error at most $BLh$.
For the displayed series the increment estimate itself follows by the
binomial formula on polynomials followed by the tail bounds on the
first two derivatives. The product and quotient rules preserve the
estimate, with a positive lower bound for each denominator. Oscillatory
factors will only be integrated on bounded intervals at a time. On
each segment all these derivative bounds are obtained by replacing
the arguments by their supplied upper bounds, and replacing $e^B$ by
$3^{\lceil B\rceil}$ for $B\ge0$. This provides a primitive recursive
mesh at any requested precision. For a finite chain of $H$ arithmetic
operations on approximations, with all magnitudes bounded by $A\ge2$
and every divisor bounded below in absolute value by $1/A$, requesting
input errors $\le 2^{-j}(8A^3)^{-H-1}$ suffices for final error
$2^{-j}$. Indeed addition, multiplication, and division amplify the
larger of the two operand errors by at most $8A^3$ while the divisor
error stays below $1/(2A)$. Induction on the rational error inequality
implements the chain for variable $H$.

\begin{theorem}[Reciprocal Dirichlet series and prime product]
For the indicated $s$ let
\[
 Z_N(s)=\sum_{n=1}^N n^{-s},\qquad
 D_N(s)=\sum_{n=1}^N\mu(n)n^{-s},\qquad
 P_N(s)=\prod_{p\le N,\ p\ \mathrm{prime}}(1-p^{-s}).
\]
Then $Z_N$ names $\zeta(s)$, and $D_N$ and $P_N$ both name
$1/\zeta(s)$. Uniform tail estimates are available given the same $q$.
\end{theorem}
\begin{proof}
On the block $2^j\le n<2^{j+1}$ the sum of $n^{-\sigma}$ is at
most $2^{-j/q}$. Group the blocks in lots of $q$. Every finite tail
starting at $n\ge2^{qk}$ has absolute majorant $2q\,2^{-k}$,
and the full sum has majorant $2q$. For $N\ge2^{2qk}$ multiply
the two finite sums $Z_N D_N$. Terms with $ab\le N$ sum to $1$ by
$\sum_{d\mid m}\mu(d)=[m=1]$. Every remaining pair has
$a>\sqrt N$ or $b>\sqrt N$. Dropping signs and extending each
rectangle of summation gives
\[
                 |Z_N D_N-1|\le 8q^2 2^{-k}.
\]
Consequently the product of the two names is $1$. In particular, the
absolute bound $2q$ on the second factor separates the first factor
from zero by $1/(2q)$. Expanding $P_N$ chooses subsets of the primes
up to $N$. The signed monomial attached to a subset with product $m$
is $\mu(m)m^{-s}$. Unique factorization gives distinct products,
and every squarefree $m\le N$ occurs. The discrepancy from $D_N$
is bounded by the absolute tail after $N$, even though the largest
product can be as large as $N!$. This proves the three assertions.
\end{proof}

\begin{theorem}[Mellin formula and its finite version]
For $N\ge1$,
\[
 s\int_1^{N} M(\lfloor t\rfloor)t^{-s-1}\,dt
      =D_N(s)-M(N)N^{-s}.
\]
Its limit is $1/\zeta(s)$. At $N\ge2^{qk}$ the error from the
limiting value is at most $(2q+1)2^{-k}$.
\end{theorem}
\begin{proof}
On $[m,m+1]$ use the endpoint difference $m^{-s}-(m+1)^{-s}$ as
the integral of $s t^{-s-1}$. Isolated endpoint values have zero
contribution, as is witnessed by meshes of arbitrarily small width.
After multiplication by $M(m)$ and summation for $1\le m<N$, the
successive differences of $M$ are $\mu$, yielding the displayed
identity, also at $N=1$. The boundary term has absolute value at most
$N^{1-\sigma}\le2^{-k}$. The Dirichlet tail gives the remaining
error. The absolute integral beyond $N$ is bounded by
$|s|\int_N^\infty t^{-\sigma}dt\le q|s|N^{-1/q}$, as follows by
integrating a power over finite intervals. This also permits arbitrary
upper endpoints in the limit.
\end{proof}

\subsection{Perron inversion and the value at a jump}

Take $x=a/b>0$ rational, with $a,b$ positive integers. Define
$M_0(x)=M(\lfloor x\rfloor)$ when $x$ is nonintegral, and
$M_0(m)=M(m)-\mu(m)/2$ at a positive integer $m$. At $m$, twice
this value is $M(m-1)+M(m)$. This distinction affects the integral
displayed in the Wikipedia article even at $x=1$.

\begin{theorem}[Symmetric Perron inversion]\label{thm:perron}
For $c\ge1+1/q$ real, let $X$ be an integer bound on
$\max(1,x^c)$, for example $X=\max(1,a)^{\lceil c\rceil}$ when
$c$ is rational (a supplied integer upper bound for $c$ suffices).
For every integer $T\ge2c$,
\[
 \left|\frac{1}{2\pi i}\int_{c-iT}^{c+iT}
       \frac{x^s}{s\zeta(s)}\,ds-M_0(x)\right|
 \le \frac{6qX(a+1)+c}{T}.
\]
In particular the normalized integral has a primitive recursive
convergence modulus. The same conclusion holds on named inputs
given a known floor and positive separation from the adjacent integers,
with the corresponding lower bound on $|\log(x/n)|$.
\end{theorem}
\begin{proof}
Put $y=\log(x/n)$ and write $K_T(y)$ for the normalized integral
of $e^{ys}/s$ along this vertical segment. When $y>0$ close the
rectangle with the left edge at $c-T$, traversing it counterclockwise.
When $y<0$ close on the right at $c+T$. The normalized integral on
the closed rectangle is respectively $1$ or $0$. To verify this
directly, subtract $1/s$ from the integrand. The remainder has
primitive
\[
           \sum_{j\ge1} \frac{y^j s^j}{j\,j!}.
\]
Its convergence and the convergence of its differentiated series on
the rectangle follow from the exponential tail estimate, so endpoint
differences cancel on the four sides. For $ds/s$ the real part is
$d\log|s|$. The imaginary part is the increment of the angle.
Explicitly in a right or left half-plane chart this differential is
$d\arctan(v/u)$ for $s=u+iv$, and in a top or bottom chart it is
$-d\arctan(u/v)$. Subdivide edges at the axes and use the continuous
angle in each quadrant. Adjacent angle values agree at the axes on
the same turn. A contour wholly in the right half-plane lies in one
chart and the endpoint increments telescope to zero. A
counterclockwise rectangle surrounding $0$ runs once through the
four successive axis directions, which advances the angle by
$4(\pi/2)=2\pi$. These chart computations are rationally effective:
$\arctan t=\int_0^t(1+u^2)^{-1}du$ and $\pi=4\arctan 1$.
The quarter-turn convention follows from
$\arctan t+\arctan(1/t)=2\arctan 1$ for $t>0$, by differentiation.

Each horizontal segment contributes absolute integral at most
$e^{cy}/(T|y|)$ since $|s|\ge T$ and the numerator decays as
$e^{-|y|u}$ at distance $u$ from the original line. The opposite
vertical edge has length $2T$, denominator at least $T/2$, and
numerator $e^{cy}e^{-T|y|}$. Since $e^{-v}\le1/v$ for $v>0$,
division by $2\pi$ gives the convenient bound
\[
       |K_T(y)-[y>0]|\le \frac{3e^{cy}}{T|y|}\quad(y\ne0).
\]
At $y=0$ symmetry on the vertical segment gives
$K_T(0)=\arctan(T/c)/\pi$, whose distance from $1/2$ is
$\arctan(c/T)/\pi\le c/T$.

If $bn<a$, then $\log(a/bn)\ge1-bn/a\ge1/a$. If $bn>a$,
then $\log(bn/a)\ge1-a/(bn)\ge1/(a+1)$. Here
$\log t\ge1-1/t$ for $t\ge1$ follows by integrating $1/u$.
Thus all nonexceptional logarithms are uniformly separated.
For fixed $T$ integrate the partial reciprocal Dirichlet sums
termwise. Their uniform tail bound on the segment gives an integral
error bounded by $2T X(2q)2^{-k}/(2\pi c)$ at cutoff
$N\ge2^{qk}$, which tends effectively to zero. The error for the
nonexceptional kernels after summation is at most
$3(a+1)X(2q)/T$. The one possible exceptional term contributes
at most $c/T$. The sum of the step values counts precisely $M_0(x)$.
\end{proof}

\section{Squarefree density and the Basel constant}\label{sec:squarefree}

Let $Q(N)=\sum_{n\le N}|\mu(n)|$, with sums over positive
integers. For each $m\ge1$,
\[
          |\mu(m)|=\sum_{d^2\mid m}\mu(d).
\]
Indeed a squarefree $d$ in this sum chooses a subset of the primes
whose squares divide $m$. The empty subset is alone when there is no
such prime. Otherwise toggling the least such prime cancels all terms
in opposite-sign pairs. Hence if $r=\lfloor\sqrt N\rfloor$, found
by a bounded scan, then
\[
                 Q(N)=\sum_{d=1}^{r}\mu(d)
                              \left\lfloor\frac{N}{d^2}\right\rfloor.
\]

\begin{theorem}[An effective squarefree count]
For $N\ge1$, $B\ge r=\lfloor\sqrt N\rfloor$, one has
\[
 \left|Q(N)-N\sum_{d=1}^B\frac{\mu(d)}{d^2}\right|
        \le r+\frac{N}{r}\le 2r+2.
\]
Consequently $|Q(N)-6N/\pi^2|\le2r+2$, proving the squarefree
asymptotic on MathWorld with a uniform error. In particular for
$N\ge2^{2j+4}$ the density differs from $6/\pi^2$ by at most
$2^{-j}$.
\end{theorem}
\begin{proof}
The $r$ floors create total error at most $r$. For $d>r$ use
$1/d^2\le 1/(d-1)-1/d$ and telescope. Finally
$N<(r+1)^2$ gives $N\le r(r+2)$. The preceding reciprocal
Dirichlet series identifies the limiting coefficient as $1/\zeta(2)$.
Dividing the explicit error by $N$ gives at most
$2/\sqrt N+2/N\le4/\sqrt N$, yielding the claimed modulus.

For completeness we identify $\zeta(2)$ within the same arithmetic.
Write $m=2l+1$, $l\ge1$. For $k=1,\ldots,l$, set $t_k=k\pi/m$.
The imaginary part of $(\cos t+i\sin t)^m$ equals $\sin(mt)$.
At $t_k$ it vanishes. Dividing its expansion by $\sin^m t_k$
gives the polynomial
\[
        \sum_{j=0}^l(-1)^j {m\choose 2j+1}Y^{l-j}
\]
at $Y=\cot^2 t_k$. The roots are distinct because $\cot t$ is
positive and strictly decreasing between $0$ and $\pi/2$.
Factoring out the $l$ monic linear factors (polynomial division, where
their distinctness permits cancellation) and comparing the next
coefficient with the leading one gives
$\sum_k \cot^2t_k={m\choose3}/m=(m-1)(m-2)/6$. On this interval,
$\sin t\le t$ by integrating $\cos t\le1$, and
$\sin t\ge t\cos t$ by integrating the derivative $t\sin t$ of
their difference. Thus $\cot^2t\le1/t^2\le1+\cot^2t$, whence
\[
 \frac{\pi^2(m-1)(m-2)}{6m^2}\le
        \sum_{k=1}^l\frac1{k^2}
          \le \frac{\pi^2(m^2-1)}{6m^2}.
\]
The two rationally approximable endpoints differ by $\pi^2l/m^2$,
and their distances to $\pi^2/6$ are at most $8/m$ since $\pi<4$.
Taking $l\ge 2^{j+4}$ suffices for error $2^{-j}$.

For the separation facts in the factorization, one can define
$\cos$ and $\sin$ as the parts of $e^{it}$. Differentiation of
$e^{i\arctan u}\sqrt{1+u^2}/(1+iu)$ shows it equals $1$ by its
value at zero. Taking $u=1$ gives $e^{i\pi/2}=i$ by squaring.
The defining integral for $\arctan 1$ gives $2\le\pi<4$.
On $0<t\le\pi/2$ the sine series gives
$\sin t\ge t-t^3/6\ge t/3$. Integrating this bound back from
$\cos(\pi/2)=0$ separates $\cos t_k$ from zero, since
$\pi/2-t_k\ge\pi/(2m)$. Also $-(\cot t)'=1/\sin^2t\ge1$,
so distinct cotangents are separated by at least $\pi/m$.
For clarity about induction strength in the polynomial step, these
bounds give $|Y_k|\le3m^2$ and $|Y_j-Y_k|\ge m^{-3}$ for $j\ne k$. These follow from
$\cos t_k\ge t_k(\pi/2-t_k)/3\ge 2/(3m^2)$ and $\sin t_k\le1$
for the latter. Take rational approximations $b_k$ to $Y_k$ within
$1/(4m^3)$. If $P$ denotes the displayed polynomial, then rational
Lagrange interpolation applied to $P(Y)-m\prod_k(Y-b_k)$ gives its
coefficient of degree $l-1$ as
\[
             \sum_k \frac{P(b_k)}{\prod_{j\ne k}(b_k-b_j)}.
\]
Indeed the interpolation polynomials take values $1$ at their own node
and $0$ at the others. Subtract their linear combination and divide
successively by the rational root factors. The right side has absolute
value at most $l(2m^3)^{l-1}\max_k|P(b_k)|$. Each residual is
bounded by $|b_k-Y_k|\,m^2 2^m(3m^2+1)^m$ by the difference of
powers formula. Thus the coefficient comparison above is a family of
finite rational inequalities with primitive recursive precision loss.
The trigonometric root identities use at most $m$ angle additions,
with the same finite-chain interpretation.
\end{proof}

\paragraph{Lean targets for the analytic arguments.}
A first target with no analysis library is the squarefree error
inequality. Let $U=(B!)^2$ and $V=\sum_{d=1}^B\mu(d)(U/d^2)$,
computing the divisions in \texttt{Nat} before casting to
\texttt{Int}. Induct over the positive summation indices to prove
\[
          r\,|UQ(N)-NV|\le U(r^2+N),\qquad 1\le r\le B,
          \quad r^2\le N<(r+1)^2.
\]
Each $d^2$ divides $U$ by the factorial multiple property. The
inequality follows from the square-divisor identity, Euclidean
division remainders, and the rational telescoping estimate above.
It gives a single executable test with natural/integer inputs.

For the transcendental targets use pairs of rationals as complex
centers, with a rational nonnegative error radius. Squared distance
between centers is rational, so containment and error tests require
only integer cross multiplication. An approximation routine accepts
a natural precision, returns a center and radius at most $2^{-j}$,
and has a proof that any two returned disks intersect. Equality of
two such names is cross-intersection at every pair of precisions.
The primitive recursive routines needed are the truncated logarithm
and exponential with the term counts above, bounded summation,
finite products, and equally spaced left sums on a finite segment.
The squarefree cutoff uses the integer square root, the zeta cutoff
uses $2^{qk}$, and the product error uses $2^{2qk}$.
Implement the Mellin integral directly as the sum of cell endpoint
differences. For Perron implement the finite segment integral as
quadrature of a finite Dirichlet polynomial, budgeting one third
of the requested error to each of series truncation, quadrature,
and rational rounding. On a segment, the quadrature error is at
most the segment length times the mesh width times a derivative
bound, obtained term by term. For $1/\zeta$ use $D_N$. The product
lemma proves the naming relation. Prove rectangle cancellation for
polynomials by endpoint subtraction, extend to the displayed
primitive by its known series tail, and prove the $1/s$ rectangle
value from the chart derivatives. Apply Theorem~\ref{thm:perron}
with $T\ge 2^j(6qX(a+1)+c)$ for the final limiting step. Each target
reduces, after supplying the precision, to bounded loops and rational
inequalities. The name-level Lean predicates collect these families.
